% ArabiziKit: an open, hybrid, benchmark-driven system for transliterating
% Arabizi into Arabic script. Compile with XeLaTeX:
%   xelatex paper.tex
% (Amiri and Noto Naskh Arabic are available on Overleaf and TeX Live full.)

\documentclass[11pt]{article}

\usepackage[margin=1in]{geometry}
\usepackage{booktabs}
\usepackage{microtype}
\usepackage{url}
\usepackage{xcolor}
\usepackage{fontspec}
\usepackage[hidelinks]{hyperref}
\setlength{\emergencystretch}{3em}

% Arabic text in examples. Tries Amiri, then Noto Naskh Arabic, then
% DejaVu Sans (ships with TeX Live), so the document compiles even
% without a dedicated Arabic font; glyphs then fall back to whatever
% supports Arabic.
\IfFontExistsTF{Amiri}
  {\newfontfamily{\arabicfont}{Amiri}[Script=Arabic]}
  {\IfFontExistsTF{Noto Naskh Arabic}
     {\newfontfamily{\arabicfont}{Noto Naskh Arabic}[Script=Arabic]}
     {\newfontfamily{\arabicfont}{DejaVu Sans}[Script=Arabic]}}
\newcommand{\ar}[1]{{\arabicfont #1}}
\newcommand{\az}[1]{\texttt{#1}}

\title{ArabiziKit: An Open, Hybrid, Benchmark-Driven Arabizi to Arabic-Script Transliteration System}

\author{Muhammed Rabeeh Mattath\\
  \small\url{https://github.com/rb2625/arabizi-kit}\\
  \small Student, American University of Ras Al Khaimah}
\date{\today}

\begin{document}
\maketitle

\begin{abstract}
Arabizi, Arabic written in Latin script with digit-substituted letters
(2 for hamza or qaf, 3 for ayn, 5 for kha), is the de facto writing
system for millions of Arabic speakers on social media, yet available
transliteration tooling is either a decade-old academic demonstration or
a component buried inside a monolithic toolkit. We present ArabiziKit, a
modern, dependency-light, open-source library that (i) combines a curated
bilingual lexicon with an interpretable rule engine that produces ranked
candidate readings rather than a single guess, (ii) performs clitic
attachment and orthographic hamza seating, handles Maghrebi conventions
(9 as qaf, ch as shin, geminates written once), and tags the input
dialect, (iii) adds a learned layer trained on a collected corpus: a
Naive Bayes dialect classifier, a word reading table, and a character
trigram language model that reranks candidates toward natural Arabic, and
(iv) ships a reproducible evaluation suite with CER, WER, exact@1, and
hit@k metrics. On our 33-sentence calibration set the hybrid pipeline
achieves exact match. On three external held-out sets with gold
references it attains 0.376 exact@1 (Egyptian, 500), 0.296
(Levantine, 389), and 0.088 (Moroccan Darija, 600) under oracle dialect
hints; the learned layer raises the pipeline's own held-out test set from
0.000 to 0.061 exact@1 and reduces CER from 0.299 to 0.226 without
exposure to test data. A corpus pipeline harvests public Arabizi,
annotates it with a free-tier LLM (inter-annotator agreement 0.171 over a
10\% double-annotated sample, which quantifies the inherent ambiguity of
the task), and produces stratified held-out splits. We release the
library, the pipeline, the benchmark, and the browser and npm
distributions.
\end{abstract}

\section{Introduction}
\label{sec:intro}

Arabizi is among the most common writing systems for Arabic on social
media, yet it remains underserved by current NLP tooling. Orthographic
conventions differ sharply by region: 2 is a hamza in Egypt and the
Levant but qaf in the Gulf and Iraq, 8 is qaf in the Egyptian convention
but heh or ghayn in Maghrebi writing, and short vowels are usually
omitted.

Transliteration into Arabic script benefits downstream applications,
including sentiment analysis, dialect identification, machine
translation, and LLM-based systems, which perform markedly better on
Arabic-script input than on Romanized text.

Our contributions are:

\begin{enumerate}
  \item A dependency-free, MIT-licensed Python library plus CLI, browser
        demo, and npm package, providing immediate reproducibility with
        no infrastructure requirements.
  \item A ranked-candidate design that treats ambiguity as first-class
        output instead of hiding it, consumed by an optional LLM mode and
        a learned reranker.
  \item A reproducible benchmark with three external held-out sets and a
        pipeline-produced held-out split with inter-annotator agreement,
        all regenerated by one command.
\end{enumerate}

\section{Background and related work}
\label{sec:related}

Arabizi orthography mixes letter and digit conventions by dialect
(2, 3, 5, 7, 9; g/j alternation; short-vowel elision; French-influenced
digraphs in Maghrebi such as ch and the 9-as-qaf convention).

Several prior systems address Arabizi and its dialectal variants.
Habash et al.~\cite{habash2012} proposed a hybrid approach for written
Egyptian colloquial; COLABA (Al-Badrashiny et al.~\cite{albadrashiny2014})
built a transliteration system for dialectal Arabic social media
together with the Jordanian Arabizi-Transliteration Corpus; Eskander
et al.~\cite{eskander2013} introduced CODAFY, which converts spontaneous
Egyptian colloquial writing into conventional orthography; CAMeL Tools
(Obeid et al.~\cite{obeid2020}) offers an open toolkit with
transliteration support; and MADAR (Bouamor et al.~\cite{bouamor2018})
supplies dialect identification resources.

To our knowledge, no maintained, open-source, modern system offers ranked
candidate output, a living benchmark, and reproducible held-out
evaluations.

\section{Method}
\label{sec:method}

\subsection{Lexicon layer}
\label{sec:lexicon}

A curated bilingual lexicon of roughly 180 entries, dialect-tagged, with a
Maghrebi block added in v0.3. Exact matches win outright and carry
dialect evidence; there is no ambiguity for known words.

\subsection{Rule engine}
\label{sec:rules}

A longest-match phoneme scan (digraphs before digit codes before single
letters) with context rules:

\begin{itemize}
  \item word-final 2 after a vowel to qaf (\az{tare2} to
        \ar{طريق}),
  \item word-initial 2 before a consonant to \{qaf, alef\}
        (\az{2alb} to \ar{قلب} vs \az{2ktob} to \ar{اكتب}),
  \item word-final a to \{taa-marbuta, alef\} (\az{7elwa} to
        \ar{حلوة} vs \az{hada} to \ar{هذا}),
  \item ay/ey to \{alif+ya, ya\} (\az{3ayza} to \ar{عايزة} vs
        \az{3alayk} to \ar{عليك}),
  \item short vowels e/o/u elided by default (\az{3emel} to
        \ar{عمل}).
\end{itemize}

\subsection{Dialect conventions}
\label{sec:dialect-conventions}

A dialect hint re-orders ambiguous readings. The Maghrebi convention
reads 9 as qaf, 8 as heh, ch as shin, writes doubled consonants once
(\az{bzzaf} to \ar{بزاف}), writes the assimilated definite article as
word-initial doubling (\az{jjaya} to \ar{الجاية}), and marks emphatic
readings with capital letters (\az{T} to taa, \az{S} to sad;
\az{Tab3an} to \ar{طبعا}). Without a hint the Egyptian convention is the
default and the ambiguity stays in the candidate list.

\subsection{Candidate generation and ranking}
\label{sec:candidates}

A bounded Cartesian product of per-position options, with
orthographic-plausibility penalties on bad bigrams and stable
tie-breaking. Output is the ranked top-k.

\subsection{Post-processing}
\label{sec:post}

Hamza seating with carrier merging (waw+hamza to waw-hamza, ya+hamza to
ya-hamza, alef+hamza to alef-madda or alef) and clitic attachment
(\az{el/al/il/l} to the definite article, \az{w} to the conjunction,
\az{3al} to on-the).

\subsection{Learned layer}
\label{sec:learned}

Three small, dependency-free components trained from the corpus (the
calibration set plus the pipeline train and dev splits; held-out and
external sets are excluded):

\begin{itemize}
  \item a Naive Bayes dialect classifier over word tokens and Arabizi
        code markers, trained on a class-balanced sample so large public
        sources do not swamp minority dialects. It supplies the dialect
        hint automatically, but only for dialects whose conventions
        change the engine's default readings (Maghrebi, Gulf);
  \item a word reading table mapping Arabizi words to observed Arabic
        renderings with frequencies, aligned with article-aware handling
        (\az{el etnein} to \ar{الاثنين}); known words emit their
        observed readings instead of the letter rules;
  \item a Laplace-smoothed character trigram language model over the
        corpus references that reranks each word's candidates toward
        natural Arabic, with the rerank weight tuned on the dev split.
\end{itemize}

\subsection{LLM-assisted mode}
\label{sec:llm}

For genuinely ambiguous input, the top-k candidates are handed to an LLM
that picks the intended reading and tags the dialect. The client is
provider-agnostic and uses free tiers by default (e.g., Groq's
\az{openai/gpt-oss-120b}), with OpenAI, Gemini, local Ollama, and
Anthropic as fallbacks. The same client powers corpus annotation.

\section{Data and benchmark}
\label{sec:data}

\subsection{Calibration set}
\label{sec:calibration}

33 hand-curated sentences across five dialect groups, sharing
distribution with the seed lexicon, so scores are an upper bound on the
hybrid pipeline (1.000 exact, 1.000 hit@3, CER 0.000).

\subsection{Corpus pipeline}
\label{sec:pipeline}

We initially targeted Reddit; however, anonymous access is blocked from
many networks and app creation is gated, so the pipeline harvests from
public Hugging Face datasets instead. The filter keeps Latin-script
sentences that score as Arabizi, strips URLs, mentions, hashtags, and
emoji, and splits posts into sentences. Annotation runs on a free-tier
LLM with retries, rate-limit pacing, resumable batch writes, and a
deterministic 10\% double-annotated sample for inter-annotator agreement.

\subsection{External held-out sets}
\label{sec:external}

Three parallel Arabizi/Arabic datasets with gold references imported
from Hugging Face, never used in training: Egyptian (500,
\url{https://huggingface.co/datasets/arbml/Arabizi_Transliteration}),
Levantine (389, \url{https://huggingface.co/datasets/akhanafer/arabic-to-arabizi}),
and Moroccan Darija (600, \url{https://huggingface.co/datasets/elkababi2/Darija-Text-Ar-Arabizi}).

\subsection{Pipeline-produced split}
\label{sec:split}

355 sentences annotated by the LLM; after dropping empty references, 323
rows stratify by dialect into train (226), dev (48), and test (49).

\subsection{Benchmark protocol}
\label{sec:protocol}

Per-sentence CER on normalized orthography (alef variants, taa-marbuta,
hamza seats, and diacritics removed), WER, exact@1, and hit@k,
aggregated overall and per dialect. \az{arabizikit eval} reproduces every
number, and \az{arabizikit eval --model} runs the learned layer.
Dialect-conditioned evaluation passes each row's dialect as the hint for
the oracle baseline.

\section{Experiments}
\label{sec:experiments}

All numbers are reproduced by the released code. Modes: rules without a
dialect hint, rules with the oracle hint, and the learned layer (no
oracle; the classifier supplies effective hints).

\begin{table}[t]
\centering
\small
\caption{Exact@1 / Hit@3 / CER on held-out data.}
\label{tab:main}
\begin{tabular}{lccc ccc ccc}
\toprule
 & \multicolumn{3}{c}{rules, no hint} & \multicolumn{3}{c}{rules, oracle hint} & \multicolumn{3}{c}{learned layer} \\
\cmidrule(lr){2-4}\cmidrule(lr){5-7}\cmidrule(lr){8-10}
Set & Exact & Hit@3 & CER & Exact & Hit@3 & CER & Exact & Hit@3 & CER \\
\midrule
Pipeline dev (48)    & 0.021 & 0.021 & 0.291 & 0.021 & 0.021 & 0.291 & \textbf{0.354} & \textbf{0.542} & \textbf{0.097} \\
Pipeline test (49)   & 0.000 & 0.000 & 0.299 & 0.000 & 0.000 & 0.299 & \textbf{0.061} & \textbf{0.102} & \textbf{0.226} \\
Egyptian (500)       & \textbf{0.376} & \textbf{0.526} & \textbf{0.236} & 0.376 & 0.526 & 0.236 & 0.348 & 0.502 & 0.248 \\
Levantine (389)      & \textbf{0.296} & \textbf{0.429} & \textbf{0.076} & 0.296 & 0.429 & 0.076 & 0.270 & 0.360 & 0.096 \\
Moroccan Darija (600) & 0.035 & 0.057 & 0.235 & \textbf{0.088} & \textbf{0.115} & \textbf{0.180} & 0.053 & 0.080 & 0.207 \\
\bottomrule
\end{tabular}
\end{table}

Analysis. The rules generalize to Egyptian and Levantine with zero
training: roughly a third of inputs are transliterated exactly, and most
errors fall within one or two characters. The Maghrebi conventions
(Section~\ref{sec:dialect-conventions}) raise Darija exact@1 from 0.013
to 0.088 under the oracle hint and cut CER from 0.290 to 0.180, the
largest single rule-level gain.

The pipeline test set scores 0.000 exact under the rules: the LLM
references are natural renderings (\az{mettabbal} to \ar{متقبل}) while
the rules emit literal letter mappings (\ar{متتاببال}). The learned
layer closes part of that gap (0.061 exact@1, CER reduced by a quarter)
entirely from training data, and on the dev split raises exact@1 from
0.021 to 0.354, a 17-fold improvement.

On the two external sets whose gold orthography differs from the
corpus's annotation style, the learned layer shows slightly lower
exact@1, and it lifts Darija above the no-hint baseline by automatically
selecting the Maghrebi convention.

\begin{table}[t]
\centering
\small
\caption{Dialect classifier agreement on held-out rows (effective hints
only; confidence threshold 0.4).}
\label{tab:classifier}
\begin{tabular}{lccc}
\toprule
Set & hinted & correct & accuracy \\
\midrule
Moroccan Darija & 580 / 600 & 547 & 94\% \\
Egyptian        & 345 / 500 & 237 & 69\% \\
Levantine       & 364 / 389 & 36  & 10\% \\
\bottomrule
\end{tabular}
\end{table}

The Levantine gap has a clear cause: the external Levantine set uses a
formal transliteration orthography (\az{al2sbou3}, \az{mshrou3}) unlike
the casual Arabizi in any training source, a documented coverage
limitation rather than a modeling failure. Note that only Maghrebi and
Gulf hints change the engine's readings, so a wrong Egyptian or Levantine
hint is harmless.

Inter-annotator agreement over the double-annotated 10\% sample is 0.171
exact on normalized Arabic. This agreement quantifies the inherent
ambiguity of the task: the same model rarely produces the identical
rendering twice, and both renderings are typically valid.

The error taxonomy drawn from the evaluation rows comprises short-vowel
elision (elided versus written), the dialect-dependent 2/qaf split,
taa-marbuta versus alef, hamza seating, proper nouns, English
code-switching, and the French-influenced Maghrebi digraphs. Each
category is represented in the per-row output of \az{arabizikit eval}.

\section{Limitations and ethics}
\label{sec:limits}

\begin{itemize}
  \item The seed set is a calibration set; the external sets measure
        generalization but their gold style differs from the corpus's
        annotation style, so neither fully represents live social-media
        Arabizi.
  \item The rule engine is not learned end to end, and dialect coverage
        is asymmetric: Maghrebi and Levantine rule conventions are thin
        outside the curated additions, and the classifier's coverage
        follows its training data.
  \item Arabic script is written without diacritics here by design; we
        do not model vocalization.
  \item Ethical note: transliteration is a rendering task; we do not
        claim sentiment or intent. Corpus collection uses public datasets
        with their terms respected and personal data is not
        redistributed.
\end{itemize}

\section{Conclusion and future work}
\label{sec:conclusion}

We presented ArabiziKit, an open, hybrid transliteration system with
ranked ambiguity handling, dialect-aware conventions, a learned layer,
an optional LLM mode, a corpus pipeline with inter-annotator agreement,
and fully reproducible evaluations.

Future work includes a learned reranker over the full candidate list
trained on more per-dialect data, extended rule coverage for Levantine
and Gulf conventions, larger annotated corpora per dialect, and
publication of the trained model and corpus artifacts on the model and
data hubs.

\begin{thebibliography}{9}

\bibitem{habash2012}
Nizar Habash, Ramy Eskander, and Abdelati Hawwari.
\newblock A hybrid approach for converting written Egyptian colloquial
dialect into diacritized Arabic.
\newblock In \emph{Proceedings of the 6th International Conference on
Informatics and Systems (INFOS)}, Cairo, Egypt, 2012.

\bibitem{albadrashiny2014}
Mohamed Al-Badrashiny, Ramy Eskander, Nizar Habash, and Owen Rambow.
\newblock Automatic Transliteration of Romanized Dialectal Arabic.
\newblock In \emph{Proceedings of the Eighteenth Conference on
Computational Natural Language Learning (CoNLL-2014)}, pages 30--38,
Ann Arbor, Michigan, USA, 2014.

\bibitem{eskander2013}
Ramy Eskander, Nizar Habash, Owen Rambow, and Nadi Tomeh.
\newblock Processing Spontaneous Orthography.
\newblock In \emph{Proceedings of the 2013 Conference of the North
American Chapter of the Association for Computational Linguistics:
Human Language Technologies (NAACL-HLT)}, pages 585--595, Atlanta,
Georgia, USA, 2013.

\bibitem{obeid2020}
Ossama Obeid, Nasser Zalmout, Salam Khalifa, Dima Taji, Mai Oudah,
Bashar Alhafni, Go Inoue, Fadhl Eryani, Alexander Erdmann, and Nizar
Habash.
\newblock CAMeL Tools: An open source Python toolkit for Arabic natural
language processing.
\newblock In \emph{Proceedings of the 12th Language Resources and
Evaluation Conference (LREC)}, pages 7022--7032, Marseille, France,
2020.

\bibitem{bouamor2018}
Houda Bouamor, Nizar Habash, Mohammad Salameh, Wajdi Zaghouani, Owen
Rambow, Dana Abdulrahim, Ossama Obeid, Salam Khalifa, Fadhl Eryani,
Alexander Erdmann, and Kemal Oflazer.
\newblock The MADAR Arabic Dialect Corpus and Lexicon.
\newblock In \emph{Proceedings of the 11th International Conference on
Language Resources and Evaluation (LREC 2018)}, Miyazaki, Japan, 2018.

\end{thebibliography}

\appendix
\section{Reproducibility}
\label{app:repro}

The library is installed with \az{pip install arabizikit} (release v1.0.0
on PyPI). Every number in this paper is reproduced with:

\begin{verbatim}
arabizikit eval           # rules, all sets
arabizikit eval --model   # learned layer
\end{verbatim}

The corpus pipeline, from harvest through annotation to stratified
splits, is documented in \az{RELEASE.md} and the repository README, and
the trained model is rebuilt locally with \az{arabizikit model train}.

\end{document}
